\section{Boundary cases, controls, and limitations}
\label{supp:boundary}
This section collects boundary cases, secondary controls, and limitations
that accompany the main-paper results. They are reported in full here so
that the main paper can present each finding at the level it is
established; no result was removed.

\subsection{CIFAR-10 covariance start}
\label{supp:boundary-cifar}
The repair improves agreement with real data while leaving a large
teacher-relative deficit. Nine of ten teacher-defined DINO regions
remain depleted, and conditional spread remains contracted. The largest
region holds 35.35\% of teacher images, 97.30\% of standard two-call
outputs, and 91.29\% after the covariance start (Supp.~Sec.~C.6). The
fraction within the real-training 1--99\% frequency range rises from
0.72\% to 24.7\%, compared with 97.2\% for DDIM200 (Supp.~Secs.~F.6 and G.1).
Against CIFAR-10.1v6, precision falls from 0.9756 to 0.9580 while KID and
coverage improve. All four prespecified KID contrasts retain their
direction against this reference, including the adverse output-map
comparison.

\subsection{Intervention location and solver choice}
\label{supp:boundary-solver}
The original four- and eight-call frozen comparisons retain the practical
prediction transforms and clipping. In the matched-grid controls,
four-call DDIM has KID 0.1132. Covariance-aware sampling does not improve
this checkpoint within the four/eight-call ceilings when both forward and
Jacobian--vector-product calls are counted; Sec.~F retains those controls
and each declared cost convention. On the EDM checkpoint, Heun's fewer,
larger steps give worse KID than Euler at three and seven calls. At seven
calls, Heun achieves DINO precision of 0.906, while only 24.7\% of its
images fall within the real-training frequency range, the same fraction
as the two-call covariance start by coincidence. This mismatch shows the
value of the complementary frequency check (Sec.~G.1).

\subsection{MNIST restoration and continuation}
\label{supp:boundary-mnist}
Ordinary restoration lowers precision by 0.01407 versus continued
withholding, while the declared five-percentage-point fidelity margins
pass against the full-data control. In the follow-up, the ordinary
continuation's classifier acceptance is 87.4\%, compared with 95.5\% on
real data and 87.3\% for the full-data control, so rejection remains
elevated relative to real data. All-class equivalence and coverage
superiority remain unresolved.

Reference refinement supplies a further control before attributing
regional imbalance to training. On CIFAR-10, the accepted frog-region mass
changes from 0.138 under DPM++25 to 0.103 under DDIM200, against 0.0943 in
observed data. A separately fitted DINO classifier reproduces the
direction. A discrepancy that shrinks with the reference sampler is
therefore not, by itself, evidence of persistent data imbalance, which is
why persistence across qualified sampler probes is part of the training
comparison.

\subsection{Crossed probes}
\label{supp:boundary-crossed}
On the 18 sampler-test cases, the 25-call sampler improved class-mass
error, coverage, and classifier acceptance with a precision loss relative
to the four-call sampler. This comparison increases inference cost from
two or four calls to 25. It also leaves open whether training could
further help; targeted checkpoints, for example, still overrepresent
digit~3, which the depletion rule does not test. On the six cases starting
from the original withheld state, the mean class-mass error difference
between restoration with the resolved sampler and the resolved sampler
alone is $-0.00331$ [$-0.00683,0.00020$], an ordinary paired 95\%
interval across three training seeds after averaging the two short
settings within each seed, so the additional class-mass benefit is
unresolved.

\subsection{Text-to-image generation}
\label{supp:boundary-t2i}
The reversal analysis was specified after the outcomes were available,
and the intervals in the main-paper text-to-image figure are descriptive;
Sec.~I retains the separate simultaneous 20-contrast checks. For DDIM4,
the PickScore preference proxy is 0.247 [0.227,0.268], below the frozen
0.45 tolerance. The DMD2 ratio against its matched CFG1.75 teacher,
0.901 [0.873,0.929], is below equality but unresolved at the 0.9
contraction margin. LCM4's change from 1.101 against CFG7.5 to 0.695
against CFG1 reflects external-guidance sensitivity; CFG1 is not a
verified training-matched teacher, and the comparison fails CLIP support
compatibility.

\subsection{Diagnostic resolution}
\label{supp:boundary-resolution}
The smallest Gaussian null yields 4/200 expansion flags from the
approximate spread interval (Sec.~C). For a 1\%-mass region at 5,000
samples per set, increasing the comparison family from 10 to 440 reduces
detection from 99.9\% to 73.95\%.

\subsection{Sampler-selection policy}
\label{supp:boundary-policy}
At the same inference cost, permuting the audit's text-to-image choices
across prompts leaves the coverage and learned-preference differences
unresolved in both families. The Sana audit scores below the 20-call
fallback on ImageReward by 0.0261, with adjusted interval
[$-0.0461,-0.0093$] (Secs.~J--L).

\subsection{Scope and limitations}
\label{supp:boundary-limits}
The initialization result is shown on one pretrained CIFAR-10 checkpoint,
and the digit findings rest on classifier labels whose rejection rate on
generated images still exceeds that of real data. The spread intervals
are approximate and fall outside the exact error guarantee, and the
Student-$t$ intervals over panels or seeds capture sampling variability
for a fixed checkpoint or training recipe, not variation across
independently pretrained models. The matched-cost prompt shuffle leaves
the advantage of per-prompt selection unresolved.
